%\documentclass[a5paper,10pt,pdftex,DIVcalc,ngerman,headsepline,liststotoc,tablecaptionabove]{scrbook}
\documentclass[a5paper,10pt,DIVcalc,ngerman,headsepline,liststotoc,tablecaptionabove]{scrbook}

\usepackage[latin1]{inputenc}
\usepackage[T1]{fontenc}
\usepackage{babel}

\usepackage[format=hang,textfont=sl,labelfont=bf,width=.85\textwidth]{caption}

% -> �ndert KOMA-Formatierungen der captions

\setkomafont{caption}{\slshape}
\setkomafont{captionlabel}{\normalfont\bfseries}
\setcapwidth{.85\textwidth}

\usepackage{float}

\usepackage[linktocpage,bookmarksnumbered,bookmarksopen]{hyperref}

\floatstyle{komaabove}
\newfloat{algorithmus}{htb}{loa}[chapter]
\floatname{algorithmus}{Algorithmus}


%\makeatletter
%  \newcommand*\captionof[1]{\def\@captype{#1}%
%    \@ifundefined{fst@#1}{}{%
%      \@nameuse{fst@#1}%
%      \@ifundefined{@float@setevery}{}{\@float@setevery{#1}}%
%      \let\caption@fs@capt\@fs@capt
%      \renewcommand\@fs@capt[2]{\egroup
%        \vskip\abovecaptionskip
%        \normalsize\caption@fs@capt{##1}{##2}%
%        \vskip\belowcaptionskip
%        \bgroup}}%
%    \caption}
%\makeatother

% neue Definition (2005-08-19) von Axel Sommerfeldt
%\makeatletter
%  \newcommand*\captionof[1]{\def\@captype{#1}%
%    \@ifundefined{fst@#1}{}{%
%      \@nameuse{fst@#1}%
%      \@ifundefined{@float@setevery}{}{\@float@setevery{#1}}%
%      \let\caption@fs@capt\@fs@capt
%      \renewcommand\@fs@capt[2]{\egroup
%        \@fs@iftopcapt
%          \vskip\belowcaptionskip
%        \else
%          \vskip\abovecaptionskip
%        \fi
%        \normalsize\caption@fs@capt{##1}{\hyper@@anchor\@currentHref{##2}}%
%        \@fs@iftopcapt
%          \vskip\abovecaptionskip
%        \else
%          \vskip\belowcaptionskip
%        \fi
%        \bgroup}}%
%    \caption}
%\makeatother


\listfiles
%%%%%%%%%%%%%%%%%%%%%%%%%%%%%%%%%%%%%%%%%%%%%%%%%%%%%%%%%%%
\begin{document}

\chapter{Bla}
bla bla bla bla bla bla bla bla bla bla bla bla bla bla bla bla bla bla bla bla bla 

\section{Abschnitt eins}

Siehe Algorithmus \ref{alg:eins} 

%\begin{table}
%	\caption{Eine Tabelle }
%	\centering
%		\begin{tabular}{|c|c|}
%			\hline
%			x&x\\
%			\hline
%		\end{tabular}
%\end{table}

\section{Abschnitt zwei}
Siehe Algorithmus \ref{alg:zwei} auf Seite \pageref{alg:zwei} und Abb. \ref{fig:eins}.
%\newpage 

\begin{algorithmus}
\hrule
\smallskip
\centering
\begin{minipage}{.9\textwidth} % etwas schmaler als Flie�text
\begin{itemize}
	\item bla1
	\item bla2
	\item bla3
\end{itemize}
\end{minipage}
\smallskip
\hrule
\caption{Ein einzelner Code bla bla bla bla bla bla bla bla bla bla bla bla bla bla bla bla bla bla bla bla bla}
\label{alg:eins}
\end{algorithmus}

\begin{figure}
	\centering
\fbox{DUMMY}\par
\caption{Einzelnes Bild bla bla bla bla bla bla bla bla bla bla bla bla bla bla bla bla bla bla bla bla bla }
\label{fig:eins}
\end{figure}

\newpage 

\begin{figure}
	\centering
\fbox{DUMMY}\par
\caption{Bild zusammen mit Code bla bla bla bla bla bla bla bla bla bla bla bla bla bla bla bla bla bla bla bla bla }
\label{fig:zwei}
\vspace{\floatsep} % zusaetzlicher Abstand zwischen zwei `floats'
\captionof{algorithmus}{Code zusammen mit Bild bla bla bla bla bla bla bla bla bla bla bla bla bla bla bla bla bla bla bla bla bla }
\label{alg:zwei}
\hrule
\smallskip
\centering
\begin{minipage}{.9\textwidth} % etwas schmaler als Flie�text
\begin{itemize}
	\item bla1
	\item bla2
	\item bla3
\end{itemize}
\end{minipage}
\smallskip
\hrule
\end{figure}

\end{document}2
	\item bla3
\end{itemize}
\end{minipage}
\smallskip
\hrule
\end{figure}

\end{document}